% Standalone problem document, generated from the adjacent README.md.
% Compile with XeLaTeX. Mathematical content is maintained in Markdown.
\documentclass[11pt,a4paper]{article}
\usepackage[fontsize=10.5pt]{scrextend}
\usepackage[margin=22mm,headheight=15pt,headsep=9mm,footskip=11mm]{geometry}
\usepackage{fontspec}
\setmainfont{texgyrepagella}[Extension=.otf,UprightFont=*-regular,BoldFont=*-bold,ItalicFont=*-italic,BoldItalicFont=*-bolditalic]
\setsansfont{texgyreheros}[Extension=.otf,UprightFont=*-regular,BoldFont=*-bold,ItalicFont=*-italic,BoldItalicFont=*-bolditalic]
\setmonofont{texgyrecursor}[Extension=.otf,UprightFont=*-regular,BoldFont=*-bold,ItalicFont=*-italic,BoldItalicFont=*-bolditalic,Scale=MatchLowercase]
\usepackage{amsmath,amssymb}
\usepackage{unicode-math}
\setmathfont{texgyrepagella-math.otf}
\usepackage{xcolor,microtype,enumitem,needspace,fancyhdr}
\usepackage{bookmark}
\definecolor{ink}{HTML}{17344B}
\definecolor{accent}{HTML}{197D80}
\definecolor{muted}{HTML}{596773}
\hypersetup{colorlinks=true,linkcolor=accent,urlcolor=accent,pdftitle={MI-17 - The
Lih--Wang permanent inequality toward the flat
matrix},pdfauthor={Open Problems in Numerical Linear Algebra}}
\urlstyle{same}
\setlength{\parindent}{0pt}
\setlength{\parskip}{5pt plus 1pt minus 1pt}
\linespread{1.04}
\setlength{\emergencystretch}{3em}
\widowpenalty=10000
\clubpenalty=10000
\displaywidowpenalty=10000
\allowdisplaybreaks[1]
\setlist{nosep,leftmargin=1.5em}
\providecommand{\tightlist}{\setlength{\itemsep}{0pt}\setlength{\parskip}{0pt}}
\providecommand{\pandocbounded}[1]{#1}
\pagestyle{fancy}
\fancyhf{}
\fancyhead[L]{\sffamily\footnotesize\color{muted}OPEN PROBLEMS IN NUMERICAL LINEAR ALGEBRA}
\fancyhead[R]{\sffamily\bfseries\footnotesize\color{accent}MI-17}
\fancyfoot[L]{\parbox[b]{0.9\textwidth}{\sffamily\fontsize{8}{10}\selectfont\color{muted}Editorial ratings; a bounded search cannot certify that a problem remains unsolved.\\Literature check: 2026-09-08}}
\fancyfoot[R]{\sffamily\footnotesize\color{muted}\thepage}
\renewcommand{\headrulewidth}{0.3pt}
\renewcommand{\headrule}{\hbox to\headwidth{\color{accent}\leaders\hrule height \headrulewidth\hfill}}
\makeatletter
\renewcommand{\section}{\Needspace{5\baselineskip}\@startsection{section}{1}{0pt}{12pt plus 2pt minus 2pt}{4pt}{\sffamily\bfseries\large\color{ink}}}
\renewcommand{\subsection}{\Needspace{4\baselineskip}\@startsection{subsection}{2}{0pt}{9pt plus 1pt minus 1pt}{3pt}{\sffamily\bfseries\normalsize\color{ink}}}
\makeatother
\setcounter{secnumdepth}{0}
\begin{document}
{\sffamily\small\color{muted}Matrix inequalities and norms\par}
\vspace{3pt}
{\raggedright\hyphenpenalty=10000\exhyphenpenalty=10000\sffamily\bfseries\fontsize{21}{24}\selectfont\color{ink}The
Lih--Wang permanent inequality toward the flat matrix\par}
\vspace{7pt}
{\sffamily\small\textbf{Difficulty:} extreme\quad\textbf{Importance:} interesting
to specialist\par}
\vspace{4pt}
{\color{accent}\hrule height 0.8pt}
\vspace{6pt}
\textbf{Provenance:} explicit conjecture

\subsection{Problem statement}\label{problem-statement}

For every integer \(n\ge3\), let \(A=(a_{ij})\in\mathbb R^{n\times n}\)
be doubly stochastic: \(a_{ij}\ge0\) and every row and every column has
sum one. Put \(J_n=\mathbf1\mathbf1^T/n\), with
\(\mathbf1\in\mathbb R^n\) the all-ones vector. Is

\[\operatorname{per}\bigl(tJ_n+(1-t)A\bigr)
\le t\operatorname{per}(J_n)+(1-t)\operatorname{per}(A)
\qquad\text{for every }t\in[1/2,1]?\]

Here
\(\operatorname{per}(C)=\sum_{\sigma\in S_n}\prod_{i=1}^n c_{i,\sigma(i)}\),
so \(\operatorname{per}(J_n)=n!/n^n\). Neither symmetry nor positive
semidefiniteness is assumed for \(A\).

\subsection{Relevance}\label{relevance}

Mixing with the flat matrix is a basic averaging operation on stochastic
matrices. The conjecture asks for a sharp bound on the permanent along a
specified part of each such segment.

\subsection{References}\label{references}

\begin{enumerate}
\def\labelenumi{\arabic{enumi}.}
\tightlist
\item
  K.-W. Lih and E. T. H. Wang, \emph{A convexity inequality on the
  permanent of doubly stochastic matrices}, Congressus Numerantium 36
  (1982), 189--198. Original statement and order-three result;
  bibliographic attribution and formulation reproduced in reference 2,
  §1.
\item
  D. K. Udayan and K. Somasundaram, \emph{Lih Wang's and Dittert's
  conjectures on permanents}, Special Matrices 12 (2024), 20240006, §1,
  equation (3), Theorems 2.1--2.2.
  \href{https://www.degruyterbrill.com/document/doi/10.1515/spma-2024-0006/html}{Published
  full text}; \href{https://arxiv.org/abs/2312.00464}{older arXiv v1}.
\end{enumerate}

\subsection{Status check --- 2026-09-08}\label{status-check-2026-09-08}

The arXiv record remains its December 2023 v1, but the May 2024
published paper has materially narrower claims than that preprint's
abstract. Published Theorem 2.1 covers order four; Theorem 2.2 covers
only certain order-six matrices and \(t\in[0.7836,1]\). The
general-dimensional assertion survives. Searches used
\textit{Lih Wang conjecture 2026 permanent},
\textit{Lih-Wang permanent counterexample proof 2025 2026}, and the
exact journal title. No full resolution was located. The full journal
theorem statements were read; the 1982 paper was traced through their
references. The failed extension to all \(t\in[0,1]\) is a different
assertion.
\end{document}
